%% GENERATED FILE -- do not hand-edit.
%% SOURCE: experiments/database/scripts/build_candidate_datasets_table.py
\begin{landscape}
\begingroup
\footnotesize
\setlength{\tabcolsep}{4pt}
\renewcommand{\arraystretch}{1.15}
\begin{longtable}{@{}r@{\hspace{4pt}} L{86mm} L{92mm} L{62mm}@{}}
\caption[]{Sources for Table~\ref{tab:final}, in the same order. \emph{Why it is
ordered here} is the band reason for a perturb-seq or molecular-layers row, and the tier
and measurement rule for a scale row; what the row buys is the \emph{Why} column of
Table~\ref{tab:final}. \emph{Data} is where the per-record values live; entries marked
unconfirmed were not fetched live and must be checked before a loader is written. Every
link is clickable.}
\label{tab:sources}\\
\toprule
\textbf{\#} & \textbf{Dataset, citation and link} & \textbf{Why it is ordered here} & \textbf{Data} \\
\midrule
\endfirsthead
\multicolumn{4}{@{}l}{\footnotesize\emph{Table~\ref{tab:sources}, continued}}\\
\toprule
\textbf{\#} & \textbf{Dataset, citation and link} & \textbf{Why it is ordered here} & \textbf{Data} \\
\midrule
\endhead
\bottomrule
\endfoot
1 & \textbf{Boocock 2025 (single-cell eQTL mapping)}\newline Boocock J, Alexander N, Alamo Tapia L, Walter-McNeill L, Patel SP, Munugala C, Bloom JS, Kruglyak L. eLife 2025;13:e95566.\newline \href{https://doi.org/10.7554/eLife.95566}{doi:10.7554/\allowbreak eLife.95566} & Genotype recovered per cell from the cell's own reads, crossed with a per-cell transcriptome. The largest existing yeast dataset in the quadrant the campaign targets. & mirrored as boocockSinglecellEQTLMapping2025; GEO (unconfirmed) \\
\addlinespace[5pt]
2 & \textbf{Hale 2024 (CRISPRi x natural variation)}\newline Hale JJ, Matsui T, Goldstein I, et al., Kruglyak L. Nat Commun 2024;15:4234.\newline \href{https://doi.org/10.1038/s41467-024-48626-1}{doi:10.1038/\allowbreak s41467-024-48626-1} & The one released yeast library that crosses a CRISPR interference guide set with a sequenced genetic background, which is the input axis a campaign has to plan against. & SRA PRJNA986287 + Figshare \\
\addlinespace[5pt]
3 & \textbf{Jackson 2020 (TF-deletion single-cell atlas)}\newline Jackson CA, Castro DM, Saldi GA, Bonneau R, Gresham D. eLife 2020;9:e51254.\newline \href{https://doi.org/10.7554/eLife.51254}{doi:10.7554/\allowbreak eLife.51254} & Transcribed genotype barcodes with a per-cell transcriptome readout, fully crossed over 11 conditions. The design a yeast Perturb-seq is specified against. & GEO GSE125162; also eLife Source code 2 (103118\_SS\_Data.tsv.gz). Confirmed from the mirrored paper. The released matrix carries 38,225 cells; the abstract says 38,285 and the discussion 38,255, so the deposit governs. The authoritative 72-strain genotype list is Supplementary file 1 Table S2, an Excel file not held in the mirror. \\
\addlinespace[5pt]
4 & \textbf{Puddu 2019 (WGS of the deletion collection)}\newline Puddu F, et al., Jackson SP. Nature 2019;573:416-420.\newline \href{https://doi.org/10.1038/s41586-019-1549-9}{doi:10.1038/\allowbreak s41586-019-1549-9} & Whole-genome sequence for every strain in the deletion collection. A pooled campaign reads a barcode and infers a genotype; this is the measurement of how often that inference is wrong, and Jackson 2020 built its own strains rather than use the collection for related reasons. & ENA/\allowbreak SRA (accession unconfirmed) \\
\addlinespace[5pt]
5 & \textbf{N'Guessan 2025 (segregant scRNA-seq eQTL)}\newline N'Guessan A, Boocock J, Kruglyak L, Albert FW. eLife 2025. PMC12303567.\newline \href{https://pmc.ncbi.nlm.nih.gov/articles/PMC12303567/}{PMC12303567} & About 4,500 sequenced segregants profiled by single-cell RNA-seq, so genotype and transcriptome are measured in the same cell. & eLife data availability; segregant panel from Bloom lineage \\
\addlinespace[5pt]
6 & \textbf{Hackett 2020 (IDEA inducible-TF transcriptome time series)}\newline Hackett SR, Baltz EA, Coram M, Wranik BJ, Kim G, Baker A, Fan M, Hendrickson DG, Berndl M, McIsaac RS. Mol Syst Biol 2020;16:e9174.\newline \href{https://doi.org/10.15252/msb.20199174}{doi:10.15252/\allowbreak msb.20199174} & Induction rather than deletion, with a transcriptome followed over time. Jackson 2020 names transient induction as the perturbation modality most likely to produce a detectable expression response. & Published with the paper (Mol Syst Biol SI); accession unconfirmed from the mirror \\
\addlinespace[5pt]
7 & \textbf{Hu 2007 (TF deletion expression compendium)}\newline Hu Z, Killion PJ, Iyer VR. Nat Genet 2007;39:683-687.\newline \href{https://doi.org/10.1038/ng2012}{doi:10.1038/\allowbreak ng2012} & A designed single-gene perturbation crossed with a transcriptome over 269 transcription-factor deletions, which is the bulk form of what the campaign measures per cell, and the widest regulator coverage available. & GEO (unconfirmed) \\
\addlinespace[5pt]
8 & \textbf{Dong 2021 (MAGIC + SAM biosensor)}\newline Dong C, Schultz JC, Liu W, Lian J, Huang L, Xu Z, Zhao H. Metab Eng 2021;66:319-327.\newline \href{https://doi.org/10.1016/j.ymben.2021.03.005}{doi:10.1016/\allowbreak j.ymben.2021.03.005} & A genome-wide CRISPR activation, interference and deletion library sorted on a metabolite biosensor, so the label is product concentration rather than fitness. The pattern a sorted Perturb-seq would reuse. & Metab Eng SI (accession unconfirmed) \\
\addlinespace[5pt]
9 & \textbf{Momen-Roknabadi 2020 (inducible CRISPRi library)}\newline Momen-Roknabadi A, Oikonomou P, Zegans M, Tavazoie S. Commun Biol 2020;3:723.\newline \href{https://doi.org/10.1038/s42003-020-01452-9}{doi:10.1038/\allowbreak s42003-020-01452-9} & A genome-scale inducible CRISPR interference library with a released per-guide matrix. Induction control is what lets a knockdown be applied after a cell is captured rather than during outgrowth. & Addgene + Commun Biol SI \\
\addlinespace[5pt]
10 & \textbf{McGlincy 2021 (genome-scale CRISPRi library)}\newline McGlincy NJ, Meacham ZA, Reynaud KK, Muller R, Baum R, Ingolia NT. BMC Genomics 2021;22:205.\newline \href{https://doi.org/10.1186/s12864-021-07518-0}{doi:10.1186/\allowbreak s12864-021-07518-0} & The second independently designed genome-scale yeast CRISPR interference library with a released per-guide matrix, so guide design and gene effect can be separated. & ingolia-lab.org/\allowbreak yeast-crispri + Addgene + BMC SI \\
\addlinespace[5pt]
11 & \textbf{Bao 2018 (CHAnGE single-nucleotide library)}\newline Bao Z, HamediRad M, Xue P, Xiao H, Tasan I, Chao R, Liang J, Zhao H. Nat Biotechnol 2018;36:505-508.\newline \href{https://doi.org/10.1038/nbt.4132}{doi:10.1038/\allowbreak nbt.4132} & A guide-indexed genome-wide library whose edits sit below the open reading frame, so the perturbation is an allele rather than a null. & Nat Biotechnol SI (raw sequencing location unconfirmed) \\
\addlinespace[5pt]
12 & \textbf{Crook 2016 (tunable RNAi, isobutanol + 1-butanol)}\newline Crook N, Sun J, Morse N, Schmitz A, Alper HS. Appl Microbiol Biotechnol 2016;100:10005-10018. PMID 27654654.\newline \href{https://pubmed.ncbi.nlm.nih.gov/27654654/}{PMID 27654654} & A dose-graded knockdown axis. Graded rather than binary perturbation is what turns a per-cell readout into a dose-response curve. & AMB SI (DOI is an open verification item) \\
\addlinespace[5pt]
13 & \textbf{Roy 2018 (multiplexed precision editing)}\newline Roy KR, Smith JD, Vonesch SC, et al., St Onge RP. Nat Biotechnol 2018;36:512-520.\newline \href{https://doi.org/10.1038/nbt.4137}{doi:10.1038/\allowbreak nbt.4137} & Multiplexed designed edits per cell, which is combinatorial input in the sense the Perturb-seq specification requires. & Nat Biotechnol SI (accession unconfirmed) \\
\addlinespace[5pt]
14 & \textbf{Newman 2006 (protein noise, GFP library)}\newline Newman JRS, Ghaemmaghami S, Ihmels J, Breslow DK, Noble M, DeRisi JL, Weissman JS. Nature 2006;441:840-846.\newline \href{https://doi.org/10.1038/nature04785}{doi:10.1038/\allowbreak nature04785} & Per-cell protein abundance distributions across a genome-scale tagged collection. Expression noise per gene is what sets how large a perturbation effect has to be before a per-cell readout can resolve it. & Nature SI \\
\addlinespace[5pt]
15 & \textbf{Mukherjee 2021 (CRISPRi essential genes x acetic acid)}\newline Mukherjee V, Lind U, St Onge RP, Blomberg A, Nystrom T. mSystems 2021;6:e00410-21.\newline \href{https://doi.org/10.1128/mSystems.00410-21}{doi:10.1128/\allowbreak mSystems.00410-21} & Knockdown reaches the essential genes a deletion collection cannot hold, which is a third of the genome a deletion-based campaign would miss. & mSystems open-access SI \\
\addlinespace[5pt]
16 & \textbf{Lian 2017 (CRISPR-AID, beta-carotene)}\newline Lian J, HamediRad M, Hu S, Zhao H. Nat Commun 2017;8:1688.\newline \href{https://doi.org/10.1038/s41467-017-01695-x}{doi:10.1038/\allowbreak s41467-017-01695-x} & Three perturbation modalities in one cell, activation, interference and deletion, which is combinatorial input by modality rather than by gene. & Nat Commun SI \\
\addlinespace[5pt]
17 & \textbf{Hughes 2000 (compendium of expression profiles)}\newline Hughes TR, Marton MJ, Jones AR, Roberts CJ, Stoughton R, et al., Friend SH. Cell 2000;102:109-126.\newline \href{https://doi.org/10.1016/S0092-8674(00)00015-5}{doi:10.1016/\allowbreak S0092-8674(00)00015-5} & The original designed-perturbation expression compendium, and the independent measurement that makes cross-laboratory replication of the deletion transcriptome testable against the built Kemmeren 2014. & Rosetta compendium, distributed with the paper; not fetched \\
\addlinespace[5pt]
18 & \textbf{Lenstra 2011 (chromatin regulator deletion expression)}\newline Lenstra TL, Benschop JJ, Kim T, et al., Holstege FCP. Mol Cell 2011;42:536-549.\newline \href{https://doi.org/10.1016/j.molcel.2011.03.026}{doi:10.1016/\allowbreak j.molcel.2011.03.026} & The same design over chromatin regulators rather than sequence-specific factors, so the two together cover both halves of transcriptional control by deletion. & GEO (unconfirmed) \\
\addlinespace[5pt]
19 & \textbf{Sun 2013 (mRNA synthesis and decay rates across deletion strains)}\newline Sun M, Schwalb B, Pirkl N, Maier KC, Schenk A, Failmezger H, Tresch A, Cramer P. Mol Cell 2013;52:52-62.\newline \href{https://doi.org/10.1016/j.molcel.2013.09.010}{doi:10.1016/\allowbreak j.molcel.2013.09.010} & A designed single-gene perturbation crossed with a transcriptome, and the only one that resolves the level into a synthesis rate and a decay rate. Rates are what a per-cell snapshot cannot recover on its own. & GEO, series not fetched this pass \\
\addlinespace[5pt]
20 & \textbf{Jariani 2020 (yeast scRNA-seq through lag phase)}\newline Jariani A, Vermeersch L, Cerulus B, Perez-Samper G, Voordeckers K, Van Brussel T, Thienpont B, Lambrechts D, Verstrepen KJ. eLife 2020;9:e55320.\newline \href{https://doi.org/10.7554/eLife.55320}{doi:10.7554/\allowbreak eLife.55320} & Per-cell states through a carbon-source shift, so the readout is the distribution across an unsynchronized population rather than its mean. & GEO GSE144820 and GSE116246; from the collection sweep, not fetched \\
\addlinespace[5pt]
21 & \textbf{Guo 2018 (CRISPR-Cas9 tiling of essential genes)}\newline Guo X, Chavez A, Tung A, et al., Church GM, Shalem O. Nat Biotechnol 2018;36:540-546.\newline \href{https://doi.org/10.1038/nbt.4147}{doi:10.1038/\allowbreak nbt.4147} & Guide tiling produces a graded allelic series across one gene, so guide position rather than gene identity carries the effect. & Nat Biotechnol SI (unconfirmed) \\
\addlinespace[5pt]
22 & \textbf{Urbonaite 2021 (yeastDrop-Seq under drug treatment)}\newline Urbonaite G, Lee JTH, Liu P, Parras GG, Hemberg M, Acar M. Commun Biol 2021;4:1372.\newline \href{https://doi.org/10.1038/s42003-021-02895-4}{doi:10.1038/\allowbreak s42003-021-02895-4} & Single-drug arms against their combination with a per-cell readout, the environmental analog of a combinatorial genetic perturbation. & GEO GSE165686; Zenodo 10.5281/\allowbreak zenodo.4767298 and 10.5281/\allowbreak zenodo.4762526; from the collection sweep, not fetched \\
\addlinespace[5pt]
23 & \textbf{Nadal-Ribelles 2019 (sensitive yeast scRNA-seq)}\newline Nadal-Ribelles M, Islam S, Wei W, et al., Posas F, Steinmetz LM. Nat Microbiol 2019;4:683-692.\newline \href{https://doi.org/10.1038/s41564-018-0346-9}{doi:10.1038/\allowbreak s41564-018-0346-9} & The high-sensitivity yeast protocol, strand and isoform aware. Sensitivity per cell is the constraint on how small a perturbation effect a campaign can resolve. & Nat Microbiol SI /\allowbreak  GEO \\
\addlinespace[5pt]
24 & \textbf{Gasch 2017 (single-cell stress heterogeneity)}\newline Gasch AP, Yu FB, Hose J, et al., Quake SR. PLoS Biol 2017;15:e2004050.\newline \href{https://doi.org/10.1371/journal.pbio.2004050}{doi:10.1371/\allowbreak journal.pbio.2004050} & Separates intrinsic from extrinsic per-cell variation in an isogenic population under stress, which is the null a perturbation effect is measured against. & PLoS Biology SI /\allowbreak  GEO \\
\addlinespace[5pt]
25 & \textbf{Jaffe 2019 (multiplexed CRISPR interference epistasis)}\newline Jaffe M, Sherlock G, Levy SF. G3 2019;9:2843-2852.\newline \href{https://doi.org/10.1534/g3.119.400323}{doi:10.1534/\allowbreak g3.119.400323} & Two knockdowns in the same cell with a measured interaction, which is the only combinatorial CRISPR interference design located in yeast. & G3 open-access SI (unconfirmed) \\
\addlinespace[5pt]
26 & \textbf{Jackson 2023 (wild-type scRNA-seq time course for RNA kinetics)}\newline Jackson CA, Beheler-Amass M, Tjarnberg A, Suresh I, Hickey AS, Bonneau R, Gresham D. bioRxiv 2023. 10.1101/2023.09.21.558277.\newline \href{https://doi.org/10.1101/2023.09.21.558277}{doi:10.1101/\allowbreak 2023.09.21.558277} & The platform companion of Jackson 2020: 173,361 cells of one strain, which is what sets the per-cell variance a perturbation has to exceed. & GEO GSE242556 (from the scYeast Methods, Sec. 4.1.3) \\
\addlinespace[5pt]
27 & \textbf{Brettner 2024 (ultra-high-throughput yeast scRNA-seq)}\newline Brettner L, et al. Yeast 2024. doi:10.1002/yea.3927.\newline \href{https://doi.org/10.1002/yea.3927}{doi:10.1002/\allowbreak yea.3927} & Combinatorial barcoding in yeast at 96 to 384 multiplexed genotypes or environments per run. The throughput route by which a campaign becomes affordable. & Yeast (Wiley) SI \\
\addlinespace[5pt]
28 & \textbf{Mulleder 2012 (prototrophic deletion collection)}\newline Mulleder M, Capuano F, Pir P, et al., Ralser M. Nat Biotechnol 2012;30:1176-1178.\newline \href{https://doi.org/10.1038/nbt.2442}{doi:10.1038/\allowbreak nbt.2442} & The prototrophic deletion collection. Jackson 2020 used a prototrophic background because auxotrophy blocks minimal and nitrogen-limited media, so this is the strain resource a campaign in defined media needs. & EUROSCARF; Addgene kit 40276 \\
\addlinespace[5pt]
29 & \textbf{Su 2023 (single-cell transcriptomes under four stresses)}\newline Su Y, Xu C, Shea J, DeStephanis D, Su Z. BMC Genomics 2023;24:88.\newline \href{https://doi.org/10.1186/s12864-023-09184-w}{doi:10.1186/\allowbreak s12864-023-09184-w} & Full-length rather than three-prime tag counting, so isoform-level readout is on the table for a campaign that needs it. & GEO GSE201386 (from the scYeast Methods) \\
\addlinespace[5pt]
30 & \textbf{Wang 2022 (single-cell transcriptomes across replicative aging)}\newline Wang J, Sang Y, Jin S, Wang X, Azad GK, McCormick MA, Kennedy BK, Li Q, Wang J, Zhang X, et al. Aging Cell 2022;21:e13712.\newline \href{https://doi.org/10.1111/acel.13712}{doi:10.1111/\allowbreak acel.13712} & The only yeast single-cell set with age as the condition axis, which is a covariate any pooled campaign carries whether or not it measures it. & GEO GSE210032 (from the scYeast Methods) \\
\addlinespace[5pt]
31 & \textbf{Albert 2018 (eQTL in 1,012 segregants)}\newline Albert FW, Bloom JS, Siegel J, Day L, Kruglyak L. eLife 2018;7:e35471.\newline \href{https://doi.org/10.7554/eLife.35471}{doi:10.7554/\allowbreak eLife.35471} & The transcriptome half of the segregant panel that Jakobson 2025, Gerke 2017 and Eder 2020 measure protein, metabolite and flux on. & GEO /\allowbreak  eLife data availability (unconfirmed) \\
\addlinespace[5pt]
32 & \textbf{Jakobson 2025 (genome-to-proteome map)}\newline Jakobson CM, et al. Science 2025;390:eadu3198.\newline \href{https://doi.org/10.1126/science.adu3198}{doi:10.1126/\allowbreak science.adu3198} & Protein abundance across the same sequenced segregant panel Albert 2018 profiles by transcriptome, so protein and transcript quantitative trait loci are measurable on one genotype set. & Science data-availability statement (accession unconfirmed) \\
\addlinespace[5pt]
33 & \textbf{Muenzner 2024 (natural-isolate proteome)}\newline Muenzner J, et al., Ralser M. Nature 2024;630:149-157.\newline \href{https://doi.org/10.1038/s41586-024-07442-9}{doi:10.1038/\allowbreak s41586-024-07442-9} & 796 proteomes drawn from the sequenced 1,011-isolate panel that the supported Caudal 2024 transcriptomes also come from. & PRIDE PXD048219 \\
\addlinespace[5pt]
34 & \textbf{Aulakh 2025 (genome-scale ionome)}\newline Aulakh SK, et al. Cell Syst 2025;16:101319.\newline \href{https://doi.org/10.1016/j.cels.2025.101319}{doi:10.1016/\allowbreak j.cels.2025.101319} & A metabolic readout on the whole deletion collection, which is the genotype axis the supported expression compendium already covers. & Cell Systems SI (accession unconfirmed) \\
\addlinespace[5pt]
35 & \textbf{Teyssonniere 2024 (species-wide proteome against transcriptome)}\newline Teyssonniere EM, Trebulle P, Muenzner J, Loegler V, Ludwig D, Amari F, Mulleder M, Friedrich A, Hou J, Ralser M, Schacherer J. Proc Natl Acad Sci USA 2024;121:e2319211121.\newline \href{https://doi.org/10.1073/pnas.2319211121}{doi:10.1073/\allowbreak pnas.2319211121} & Proteomes for 942 isolates, the panel whose transcriptomes the built Caudal 2024 supplies, so the strain-level pairing is one to one. & PNAS SI; PMC11087752. PRIDE identifier not fetched this pass \\
\addlinespace[5pt]
36 & \textbf{Grossbach 2022 (BY x RM transcriptome, proteome and phosphoproteome)}\newline Grossbach J, Gillet L, Clement-Ziza M, Schmalohr CL, Schubert OT, Schutter M, Mawer JSP, Barnes CA, Bludau I, Weith M, Tessarz P, Graef M, Aebersold R, Beyer A. Mol Syst Biol 2022;18:e10712.\newline \href{https://doi.org/10.15252/msb.202110712}{doi:10.15252/\allowbreak msb.202110712} & Three layers on one set of 112 sequenced strains, so the transcript and the protein of a gene are observed in the same genotype and the RNA to protein question becomes a within-strain residual. & Mol Syst Biol SI + PRIDE, identifier not fetched this pass \\
\addlinespace[5pt]
37 & \textbf{Leutert 2023 (phosphoproteome x 101 conditions)}\newline Leutert M, et al., Villen J. Nat Struct Mol Biol 2023;30:1761-1773.\newline \href{https://doi.org/10.1038/s41594-023-01115-3}{doi:10.1038/\allowbreak s41594-023-01115-3} & Enzyme regulation acts faster than transcription, so a phosphosite layer over 101 conditions is what explains flux changes an expression table cannot. & PRIDE PXD034997 \\
\addlinespace[5pt]
38 & \textbf{Brauer 2008 (growth-rate-controlled chemostat transcriptome)}\newline Brauer MJ, Huttenhower C, Airoldi EM, et al., Botstein D. Mol Biol Cell 2008;19:352-367.\newline \href{https://doi.org/10.1091/mbc.e07-08-0779}{doi:10.1091/\allowbreak mbc.e07-08-0779} & Expression at controlled growth rate under each nutrient limitation, which is the condition axis Boer 2010 and Hackett 2016 measure metabolites and flux on. & GEO (unconfirmed) \\
\addlinespace[5pt]
39 & \textbf{Hackett 2016 (SIMMER multi-omic flux)}\newline Hackett SR, Zanotelli VRT, Xu W, et al., Rabinowitz JD. Science 2016;354:aaf2786.\newline \href{https://doi.org/10.1126/science.aaf2786}{doi:10.1126/\allowbreak science.aaf2786} & Flux, metabolite and transcript measured in one study under the same nutrient limitations, so the join is internal rather than across papers. & Science SI (accession unconfirmed) \\
\addlinespace[5pt]
40 & \textbf{Zhu 2014 (kinase / phosphatase lipidomics)}\newline Zhu Z, Loewen CJR. Mol Biol Cell 2014. PMC4196872.\newline \href{https://pmc.ncbi.nlm.nih.gov/articles/PMC4196872/}{PMC4196872} & A lipidome over 129 signaling deletions, most of which carry an expression profile in the supported compendium. & Mol Biol Cell SI Table S4 (Excel) \\
\addlinespace[5pt]
41 & \textbf{Gerke 2017 (urea-cycle mQTL)}\newline Gerke J, et al., Fay JC. Genetics 2017;206:2199.\newline \href{https://doi.org/10.1534/genetics.117.201107}{doi:10.1534/\allowbreak genetics.117.201107} & Metabolite quantitative trait loci on a segregant cross that also has an expression map, so a metabolite locus can be read through its transcript. & Genetics (GSA) SI \\
\addlinespace[5pt]
42 & \textbf{Ambroset 2014 (metabolite QTL)}\newline Ambroset C, et al., Fay JC. PLoS Genet 2014;10:e1004142.\newline \href{https://doi.org/10.1371/journal.pgen.1004142}{doi:10.1371/\allowbreak journal.pgen.1004142} & A 74-metabolite panel on a sequenced segregant panel, the widest metabolite vector available on a recombinant genotype axis. & PLOS Genetics SI \\
\addlinespace[5pt]
43 & \textbf{Blank 2005 (13C metabolic flux)}\newline Blank LM, Kuepfer L, Sauer U. Genome Biol 2005;6:R49.\newline \href{https://doi.org/10.1186/gb-2005-6-6-r49}{doi:10.1186/\allowbreak gb-2005-6-6-r49} & The only row measuring flux rather than a concentration, on deletion mutants that also have a transcriptome in the supported set. & Genome Biology open-access SI \\
\addlinespace[5pt]
44 & \textbf{Skelly 2013 (expression variation across isolates)}\newline Skelly DA, Merrihew GE, Riffle M, et al., MacCoss MJ, Akey JM. Genome Res 2013;23:1496-1504.\newline \href{https://doi.org/10.1101/gr.155762.113}{doi:10.1101/\allowbreak gr.155762.113} & Transcriptome and proteome on the same 22 strains, which is the only row where the cheap and expensive layers are paired within one experiment rather than joined across two. & Genome Research SI (unconfirmed) \\
\addlinespace[5pt]
45 & \textbf{Airoldi 2016 (nitrogen-limited steady-state and dynamic transcriptome)}\newline Airoldi EM, Miller D, Athanasiadou R, Brandt N, Abdul-Rahman F, Neymotin B, Hashimoto T, Bahmani T, Gresham D. Mol Biol Cell 2016;27:1383-1396.\newline \href{https://doi.org/10.1091/mbc.E14-05-1013}{doi:10.1091/\allowbreak mbc.E14-05-1013} & Expression under the exact nitrogen-limited conditions that Jackson 2020 profiles single cells in, so the deletion effect separates from the medium effect. & GEO, series not confirmed this pass \\
\addlinespace[5pt]
46 & \textbf{McManus 2014 (ribosome profiling, allele-specific)}\newline McManus CJ, May GE, Spealman P, Shteyman A. Genome Res 2014;24:422-430.\newline \href{https://doi.org/10.1101/gr.164996.113}{doi:10.1101/\allowbreak gr.164996.113} & Translation efficiency per gene alongside the mRNA abundance of that gene, which is the first of the two terms that separate protein level from transcript level. & GEO, series not fetched this pass \\
\addlinespace[5pt]
47 & \textbf{Martin-Perez 2017 (protein half-lives)}\newline Martin-Perez M, Villen J. Cell Syst 2017;5:283-294.e5.\newline \href{https://doi.org/10.1016/j.cels.2017.08.008}{doi:10.1016/\allowbreak j.cels.2017.08.008} & Protein turnover per gene, the second of those two terms, spanning three orders of magnitude and therefore not replaceable by a global constant. & Cell Systems SI; PRIDE identifier not fetched this pass \\
\addlinespace[5pt]
48 & \textbf{Boer 2010 (metabolome across nutrient limitations)}\newline Boer VM, Crutchfield CA, Bradley PH, Botstein D, Rabinowitz JD. Mol Biol Cell 2010;21:198-211.\newline \href{https://doi.org/10.1091/mbc.e09-07-0597}{doi:10.1091/\allowbreak mbc.e09-07-0597} & The metabolite half of the chemostat nutrient-limitation series whose transcriptome half is Brauer 2008, same laboratory and same conditions. & MBoC SI (unconfirmed) \\
\addlinespace[5pt]
49 & \textbf{Tengolics 2024 (domestication metabolome)}\newline Tengolics R, Szappanos B, Mulleder M, et al., Ralser M, Papp B. PNAS 2024;121:e2313354121.\newline \href{https://doi.org/10.1073/pnas.2313354121}{doi:10.1073/\allowbreak pnas.2313354121} & Metabolite levels across sequenced isolates, the panel the supported Caudal 2024 transcriptomes and Muenzner 2024 proteomes also sit on. & PNAS Dataset S11 \\
\addlinespace[5pt]
50 & \textbf{Eder 2020 (flux QTL)}\newline Eder M, Nidelet T, Sanchez I, Camarasa C, Legras JL, Dequin S. PMID 32034164.\newline \href{https://pubmed.ncbi.nlm.nih.gov/32034164/}{PMID 32034164} & Flux mapped to segregant genotypes, which is the flux counterpart of the expression and protein maps on the same panel type. & Journal SI (venue and accession unconfirmed) \\
\addlinespace[5pt]
51 & \textbf{Foss 2007 (BY x RM segregant proteome)}\newline Foss EJ, Radulovic D, Shaffer SA, Ruderfer DM, Bedalov A, Goodlett DR, Kruglyak L. Nat Genet 2007;39:1369-1375.\newline \href{https://doi.org/10.1038/ng.2007.22}{doi:10.1038/\allowbreak ng.2007.22} & The independent replicate of Grossbach 2022 on the same cross, and the first measurement that protein-level loci differ from transcript-level loci. & Nat Genet supplementary tables; not fetched \\
\addlinespace[5pt]
52 & \textbf{Yu 2021 (proteome and metabolome under nitrogen limitation)}\newline Yu R, Vorontsov E, Sihlbom C, Nielsen J. eLife 2021;10:e65722.\newline \href{https://doi.org/10.7554/eLife.65722}{doi:10.7554/\allowbreak eLife.65722} & Both layers measured in one study under nitrogen limitation, so it calibrates the protein-to-metabolite step that cross-study joins assume. & eLife data availability (unconfirmed) \\
\addlinespace[5pt]
53 & \textbf{de Boer 2020 (100M random promoters)}\newline de Boer CG, Vaishnav ED, Sadeh R, Abeyta EL, Friedman N, Regev A. Nat Biotechnol 2020;38:56-65.\newline \href{https://doi.org/10.1038/s41587-019-0315-8}{doi:10.1038/\allowbreak s41587-019-0315-8} & Scale band: tier 1, ordered by measurements (100,000,000). & GEO (de Boer 2020); DREAM 2022 release of 6,739,258 promoters \\
\addlinespace[5pt]
54 & \textbf{Vaishnav 2022 (regulatory DNA fitness landscape)}\newline Vaishnav ED, de Boer CG, Molinet J, Yassour M, Fan L, Adiconis X, Thompson DA, Levin JZ, Cubillos FA, Regev A. Nature 2022;603:455-463.\newline \href{https://doi.org/10.1038/s41586-022-04506-6}{doi:10.1038/\allowbreak s41586-022-04506-6} & Scale band: tier 1, ordered by measurements (20,000,000). & GEO per the Nature data-availability statement (unconfirmed) \\
\addlinespace[5pt]
55 & \textbf{Lee 2014 (HIP-HOP fitness signatures)}\newline Lee AY, St Onge RP, Proctor MJ, et al., Giaever G, Nislow C. Science 2014;344:208-211.\newline \href{https://doi.org/10.1126/science.1250217}{doi:10.1126/\allowbreak science.1250217} & Scale band: tier 1, ordered by measurements (13,000,000). & Science SI + Nislow/\allowbreak Giaever portal, neither of which yielded per-strain matrices; awaiting author matrices per the WS15 roadmap \\
\addlinespace[5pt]
56 & \textbf{Nguyen Ba 2022 (barcoded bulk QTL, 100k segregants)}\newline Nguyen Ba AN, Lawrence KR, Rego-Costa A, Gopalakrishnan S, Temko D, Michor F, Desai MM. eLife 2022;11:e73983.\newline \href{https://doi.org/10.7554/eLife.73983}{doi:10.7554/\allowbreak eLife.73983} & Scale band: tier 1, ordered by measurements (800,000). & eLife data availability; SRA (unconfirmed) \\
\addlinespace[5pt]
57 & \textbf{Galardini 2019 (four backgrounds x 38 conditions)}\newline Galardini M, Busby BP, Vieitez C, Dunham AS, Typas A, Beltrao P. Mol Syst Biol 2019;15:e8831.\newline \href{https://doi.org/10.15252/msb.20198831}{doi:10.15252/\allowbreak msb.20198831} & Scale band: tier 1, ordered by measurements (575,472). & GEO GSE123118 + github.com/\allowbreak mgalardini/\allowbreak 2018koyeast \\
\addlinespace[5pt]
58 & \textbf{Parsons 2006 (bioactive-compound profiling)}\newline Parsons AB, Lopez A, Givoni IE, et al., Boone C. Cell 2006;126:611-625.\newline \href{https://doi.org/10.1016/j.cell.2006.06.040}{doi:10.1016/\allowbreak j.cell.2006.06.040} & Scale band: tier 1, ordered by measurements (410,000). & Cell SI tables (accession unconfirmed) \\
\addlinespace[5pt]
59 & \textbf{Dutta 2026 (barcoded natural-isolate chemical response)}\newline Dutta A, Garin M, Loegler V, Brach G, Friedrich A, Yoshimura M, Hirano H, Osada H, Boone C, Yashiroda Y, Hou J, Schacherer J. Nat Commun 2026.\newline \href{https://doi.org/10.1038/s41467-026-73532-z}{doi:10.1038/\allowbreak s41467-026-73532-z} & Scale band: tier 1, ordered by measurements (312,000). & Nat Commun SI; isolates from the 1,011 panel (ENA) \\
\addlinespace[5pt]
60 & \textbf{Li 2016 (tRNA fitness landscape)}\newline Li C, Qian W, Maclean CJ, Zhang J. Science 2016;352:837-840.\newline \href{https://doi.org/10.1126/science.aae0568}{doi:10.1126/\allowbreak science.aae0568} & Scale band: tier 1, ordered by measurements (65,000). & Science SI (unconfirmed) \\
\addlinespace[5pt]
\end{longtable}
\endgroup
\end{landscape}
